% !PRINTED 316
\weberpage{316}

But one can give it yet another, very remarkable form.

By the relations that subsist between the functions $f(\omega), f_1(\omega), f_2(\omega)$
[§\,34, (11), (12)] we have:

\[
\frac{f(\omega)^{24} - 64}{f(\omega)^{12}} = \frac{[f_1(\omega)^8 - f_2(\omega)^8]^2}{f(\omega)^4}.
\]

If, therefore, for $w$ we introduce the new unknown

\weq{26}{y = \sqrt{w} = \frac{f_1(\omega)^8 - f_2(\omega)^8}{f(\omega)^2(w^2+5)}}

then we obtain the form:

\weq{27}{y^5 + 5y = \frac{f_1(\omega)^8 - f_2(\omega)^8}{f(\omega)^2}.}

On these formulae rests the solution of the equation of the fifth degree by
elliptic functions created by \gesperrt{Hermite} and \gesperrt{Kronecker}, to
which we have referred in the 14th chapter of Band II.

By Bd.\,I, §\,60 the general equation of the fifth degree can be reduced to the
form:

\weq{28}{z^5 + 5z = a,}

and this equation becomes identical with (27) if one sets:

\[
f_1(\omega)^8 - f_2(\omega)^8 = af^2(\omega).
\]

If one adds to this the equations:

\[
f_1^8 = f_2^8 = f^8, \quad ff_1f_2 = \sqrt{2}, % [sic] German source prints both equals signs.
\]

one obtains

\weq{29}{f^{24} - a^2f^{12} - 64 = 0.}

By this quadratic equation $f^{12}$ is determined as a function of $a$.

The roots of equation (28) are then represented by the square roots of the
expressions (22). The double sign of these square roots is explained by the
fact that two opposite values of $a$ lead to the same equation (29). But the
roots can also be determined uniquely by formula (26):

\weq{30}{z = \frac{a}{w^2+5}.}
